% p295_1; source PDF page 295; manually verified editable table.
\begin{tikzpicture}[x=1cm,y=-1cm,inner sep=0pt,outer sep=0pt]
\path[use as bounding box] (0,0) rectangle (16.000,14.100);
\draw[line width=.4pt] (0.000,0.000) rectangle (7.700,0.850);
\node[align=center,text width=7.440cm,font={\linespread{1}\fontsize{10}{13.50}\selectfont\bfseries}] at (3.850,0.425) {行为类型};
\draw[line width=.4pt] (7.700,0.000) rectangle (14.300,0.850);
\node[align=center,text width=6.340cm,font={\linespread{1}\fontsize{10}{13.50}\selectfont\bfseries}] at (11.000,0.425) {行为特征};
\draw[line width=.4pt] (14.300,0.000) rectangle (16.000,0.850);
\node[align=center,text width=1.440cm,font={\linespread{1}\fontsize{10}{13.50}\selectfont\bfseries}] at (15.150,0.425) {结论};
\draw[line width=.4pt] (0.000,0.850) rectangle (7.700,4.100);
\node[align=left,text width=7.440cm,font={\linespread{1}\fontsize{10}{13.50}\selectfont}] at (3.850,2.475) {\textbf{甲一开始就有非法占有目的}\\甲与乙在饭馆吃饭，欲非法占有乙的手机，欺骗乙：“我手机没电了，借打一下。”乙借给甲。甲假装打电话，打着打着，离开饭馆};
\draw[line width=.4pt] (7.700,0.850) rectangle (14.300,4.100);
\node[align=left,text width=6.340cm,font={\linespread{1}\fontsize{10}{13.50}\selectfont}] at (11.000,2.475) {乙的这个“借出”，没有转移占有，仍占有手机。甲将乙占有的手机转移为自己占有};
\draw[line width=.4pt] (14.300,0.850) rectangle (16.000,4.100);
\node[align=left,text width=1.440cm,font={\linespread{1}\fontsize{10}{13.50}\selectfont}] at (15.150,2.475) {盗窃罪};
\draw[line width=.4pt] (0.000,4.100) rectangle (7.700,8.100);
\node[align=left,text width=7.440cm,font={\linespread{1}\fontsize{10}{13.50}\selectfont}] at (3.850,6.100) {\textbf{甲一开始就有非法占有目的}\\甲与乙在饭馆吃饭，欲非法占有乙的手机，欺骗乙：“我手机坏了，借打3天。”乙借给甲。甲拿回家了。3天后，乙找甲还，找不到甲};
\draw[line width=.4pt] (7.700,4.100) rectangle (14.300,8.100);
\node[align=left,text width=6.340cm,font={\linespread{1}\fontsize{10}{13.50}\selectfont}] at (11.000,6.100) {乙的这个“借出”，转移了占有。甲一开始有非法占有目的，欺骗乙转移了手机的占有权。诈骗罪中的受骗人处分财物，可以只处分财物占有权，不要求处分所有权。这属于“名为借，实为骗”};
\draw[line width=.4pt] (14.300,4.100) rectangle (16.000,8.100);
\node[align=left,text width=1.440cm,font={\linespread{1}\fontsize{10}{13.50}\selectfont}] at (15.150,6.100) {诈骗罪};
\draw[line width=.4pt] (0.000,8.100) rectangle (7.700,10.700);
\node[align=left,text width=7.440cm,font={\linespread{1}\fontsize{10}{13.50}\selectfont}] at (3.850,9.400) {\textbf{甲一开始没有非法占有目的}\\甲与乙在饭馆吃饭，手机没电，借打乙的手机，打完后欲非法占有，拿着手机离开饭馆};
\draw[line width=.4pt] (7.700,8.100) rectangle (14.300,10.700);
\node[align=left,text width=6.340cm,font={\linespread{1}\fontsize{10}{13.50}\selectfont}] at (11.000,9.400) {乙的这个“借出”，没有转移占有};
\draw[line width=.4pt] (14.300,8.100) rectangle (16.000,10.700);
\node[align=left,text width=1.440cm,font={\linespread{1}\fontsize{10}{13.50}\selectfont}] at (15.150,9.400) {盗窃罪};
\draw[line width=.4pt] (0.000,10.700) rectangle (7.700,14.100);
\node[align=left,text width=7.440cm,font={\linespread{1}\fontsize{10}{13.50}\selectfont}] at (3.850,12.400) {\textbf{甲一开始没有非法占有目的}\\甲与乙在饭馆吃饭，手机坏了，向乙借用手机，答应3天后还。甲借用一天后，欲不法所有该手机。等乙索要，找不到甲};
\draw[line width=.4pt] (7.700,10.700) rectangle (14.300,14.100);
\node[align=left,text width=6.340cm,font={\linespread{1}\fontsize{10}{13.50}\selectfont}] at (11.000,12.400) {乙的这个“借出”，转移了占有权。甲属于将他人所有、自己占有的财物据为己有（行使所有权）};
\draw[line width=.4pt] (14.300,10.700) rectangle (16.000,14.100);
\node[align=left,text width=1.440cm,font={\linespread{1}\fontsize{10}{13.50}\selectfont}] at (15.150,12.400) {侵占罪};
\end{tikzpicture}
